% ============================================================================
% CHAPTER 1 TEACHER'S MANUAL
% Complete instructional guide for teaching human-in-the-loop concepts
% ============================================================================
% Source: /markdown/ch1/Ch1T_updated.md
% Note: This file is for the Instructor's Edition, not the main book
% ============================================================================

\chapter{Teacher's Manual: Chapter 1}
\label{ch:teacher-ch01}

\begin{chapterquote}{Complete instructional guide}
Teaching human-in-the-loop concepts using \emph{The Computers That Ask for Help}.
\end{chapterquote}

% ============================================================================
\section{Course Integration Guide}
% ============================================================================

\subsection{Learning Objectives}

By the end of this chapter, students will be able to:

\paragraph{Knowledge (Remembering \& Understanding).}
\begin{itemize}
    \item Define human-in-the-loop (HITL) systems and uncertainty quantification
    \item Identify examples of HITL systems in everyday technology
    \item Explain the difference between aleatoric and epistemic uncertainty
    \item Describe why intelligent systems benefit from human oversight
    \item List the five dimensions of the HITL analysis framework
\end{itemize}

\paragraph{Application \& Analysis.}
\begin{itemize}
    \item Analyze real-world systems to identify HITL intervention points
    \item Apply the five-dimension framework to evaluate HITL systems
    \item Evaluate the effectiveness of different uncertainty communication strategies
    \item Compare automated vs.\ human-in-the-loop approaches for specific scenarios
    \item Assess the trade-offs between system autonomy and human control
\end{itemize}

\paragraph{Synthesis \& Evaluation.}
\begin{itemize}
    \item Design basic HITL intervention strategies for new applications
    \item Critique existing HITL implementations using the five-dimension framework
    \item Propose improvements to human-AI collaboration interfaces
    \item Evaluate ethical implications of uncertainty-based decision systems
\end{itemize}

\subsection{Prerequisites}
\begin{itemize}
    \item \textbf{Basic Computer Science}: Introductory programming, basic algorithms
    \item \textbf{Statistics}: Understanding of probability, confidence intervals, basic hypothesis testing
    \item \textbf{Optional but Helpful}: Machine learning fundamentals, human-computer interaction basics
\end{itemize}

\subsection{Course Positioning}

\begin{description}
    \item[Early Course (Week 2--3)] Foundation for understanding human-AI collaboration
    \item[Mid-Course] Concrete examples before diving into technical implementation
    \item[Advanced Course] Quick review before focusing on implementation details
\end{description}

% ============================================================================
\section{Key Concepts for Instruction}
% ============================================================================

\subsection{The Five-Dimension Framework}

This framework is central to Chapter~1 and should be emphasized throughout:

\begin{table}[htbp]
\caption{The Five-Dimension HITL Framework}
\label{tab:five-dimension-framework}
\centering
\begin{tabular}{@{}p{3.5cm}p{4cm}p{4cm}@{}}
\toprule
\textbf{Dimension} & \textbf{Key Question} & \textbf{Example} \\
\midrule
Uncertainty Detection & Can the system recognize when it's unsure? & Netflix monitors engagement probability \\
Intervention Design & How does it ask for help? & ``Are you still watching?'' (Yes/No) \\
Timing & When does it ask? & After 3 episodes or 90 minutes \\
Stakes Calibration & Does it understand consequences? & Low stakes = can afford to interrupt \\
Feedback Integration & Does it learn from responses? & Adjusts thresholds based on user patterns \\
\bottomrule
\end{tabular}
\end{table}

\begin{technote}[Teaching Tip]
Have students apply this framework to every example. It provides a consistent analytical lens.
\end{technote}

\subsection{Contrast Pair: Netflix vs.\ Nest Thermostat}

Chapter~1 uses Netflix (good design) and Nest thermostat (problematic design) as a contrast pair:

\begin{table}[htbp]
\caption{Netflix vs.\ Nest Thermostat: Key contrast}
\label{tab:netflix-nest-contrast}
\centering
\begin{tabular}{@{}lll@{}}
\toprule
\textbf{Aspect} & \textbf{Netflix} & \textbf{Nest Thermostat} \\
\midrule
Asks for help? & Yes (``Still watching?'') & No (acts autonomously) \\
Expresses uncertainty? & Implicitly, by asking & No \\
User complaints & Minor annoyance & ``I woke up freezing'' \\
Failure mode & Interrupts unnecessarily & Confident but wrong \\
\bottomrule
\end{tabular}
\end{table}

\begin{technote}[Teaching Tip]
This contrast is powerful for showing that ``smart'' does not mean ``autonomous.''
\end{technote}

\subsection{Two Failure Modes}

Emphasize that systems can fail in two opposite ways:
\begin{enumerate}
    \item \textbf{Asking too often} $\rightarrow$ Alert fatigue (users ignore important warnings)
    \item \textbf{Never asking} $\rightarrow$ Dangerous overconfidence (GPS into lake, chatbot hallucinations)
\end{enumerate}

The Goldilocks problem: finding the ``just right'' balance.

% ============================================================================
\section{Lecture Planning Guide}
% ============================================================================

\subsection{50-Minute Lecture Structure}

\subsubsection{Opening Hook (5 minutes)}

\paragraph{Activity: Netflix Demo.}
\begin{itemize}
    \item Show actual Netflix ``Are you still watching?'' screen
    \item Poll class: ``How many of you have seen this? How many find it annoying?''
    \item Reveal: ``This isn't a bug---it's one of the smartest things Netflix does''
\end{itemize}

\subsubsection{Part 1: Pattern Recognition (12 minutes)}

\paragraph{Interactive Examples.}
\begin{itemize}
    \item \textbf{Autocorrect fails} (show funny ``ducking'' examples, discuss why they happen)
    \item \textbf{Voice assistant confusion} (play audio of Siri/Alexa saying ``I don't understand'')
    \item \textbf{Nest thermostat complaints} (show real forum posts about freezing houses)
\end{itemize}

\begin{keyconcept}[Key Teaching Point]
Some systems ask for help (Netflix, voice assistants), others don't (Nest Home/Away). What's the difference?
\end{keyconcept}

\subsubsection{Part 2: The Framework (10 minutes)}

\paragraph{Introduce the Five Dimensions.}
\begin{enumerate}
    \item Uncertainty Detection
    \item Intervention Design
    \item Timing
    \item Stakes Calibration
    \item Feedback Integration
\end{enumerate}

\textbf{Activity:} Apply framework to Netflix example on the board together.

\subsubsection{Part 3: High-Stakes Examples (15 minutes)}

\paragraph{Case Studies.}
\begin{itemize}
    \item \textbf{GPS into water} (systems that never express uncertainty)
    \item \textbf{Air Canada chatbot} (confidently wrong with legal consequences)
    \item \textbf{Medical AI} flagging uncertain diagnoses (stakes calibration done right)
    \item \textbf{Tesla Autopilot} handovers (show dashboard warnings)
\end{itemize}

\begin{trythis}
Discussion Question: ``Using the framework, what dimension failed in the Air Canada chatbot?''
\end{trythis}

\subsubsection{Part 4: Wrap-up \& Preview (8 minutes)}

\paragraph{The Two Failure Modes.}
\begin{itemize}
    \item Too much asking $\rightarrow$ Alert fatigue
    \item Too little asking $\rightarrow$ Overconfidence
\end{itemize}

\textbf{Summary:} Intelligence isn't about never being wrong---it's about knowing when you might be wrong.

\textbf{Assign:} ``Try This'' exercise (count HITL moments in next 24 hours).

\textbf{Preview:} Next class will explore how systems detect uncertainty.

\subsection{75-Minute Extended Session}

\subsubsection{Technical Deep Dive (15 minutes)}
\begin{itemize}
    \item Show actual uncertainty estimation code (from Technical Appendix)
    \item Demonstrate confidence thresholds in action
    \item Discuss calibration challenges
    \item Show the improved Nest thermostat pseudocode from the appendix
\end{itemize}

\subsubsection{Hands-On Activity (10 minutes)}
\begin{itemize}
    \item Students identify HITL moments in apps on their phones
    \item Apply the five-dimension framework to at least two examples
    \item Share findings with class
\end{itemize}

% ============================================================================
\section{Discussion Questions by Level}
% ============================================================================

\subsection{Introductory Level}

\begin{enumerate}
    \item \textbf{Recognition:} ``Give three examples of technology asking you for help. Why do you think it asks instead of just making a decision?''

    \item \textbf{Analysis:} ``Netflix could easily keep playing episodes forever. Why doesn't it? What are the costs and benefits of the `Are you still watching?' interruption?''

    \item \textbf{Comparison:} ``Compare the Netflix `Are you still watching?' prompt to the Nest thermostat's Home/Away detection. One asks for help, one doesn't. What are the consequences of each approach?''

    \item \textbf{Application:} ``If you were redesigning the Nest thermostat, when would you want it to ask the human for input versus making decisions automatically?''
\end{enumerate}

\subsection{Intermediate Level}

\begin{enumerate}
    \item \textbf{Framework Application:} ``Apply the five-dimension framework to analyze Uber's surge pricing notifications. Rate each dimension 1--5.''

    \item \textbf{Design Thinking:} ``Design the `asking for help' strategy for a new application: an AI system that helps teachers grade student essays. When should it ask for human review?''

    \item \textbf{Trade-off Analysis:} ``Self-driving cars face a dilemma: ask for human help too often and people lose trust; ask too rarely and safety suffers. How would you find the right balance?''

    \item \textbf{Ethical Considerations:} ``The Air Canada tribunal ruled that companies are responsible for what their chatbots say. What are the implications for how companies should design AI uncertainty?''

    \item \textbf{System Evaluation:} ``How would you measure whether a human-in-the-loop system is working well? What metrics matter most?''
\end{enumerate}

\subsection{Advanced Level}

\begin{enumerate}
    \item \textbf{Technical Implementation:} ``Given the uncertainty estimation methods in the technical appendix, how would you implement an optimal intervention threshold for a medical diagnosis system?''

    \item \textbf{Research Critique:} ``Read the Kendall \& Gal (2017) paper on uncertainty types. How do their findings apply to the Netflix example? What type of uncertainty is Netflix primarily addressing?''

    \item \textbf{Framework Extension:} ``The five-dimension framework covers individual interactions. How would you extend it to analyze system-level learning over time?''

    \item \textbf{Novel Applications:} ``Propose a new domain where human-in-the-loop systems could be valuable. Design the uncertainty detection and human intervention strategy using all five dimensions.''

    \item \textbf{Future Implications:} ``As AI systems become more capable, will they need human oversight more or less? Defend your position with specific examples.''
\end{enumerate}

% ============================================================================
\section{Hands-On Activities \& Exercises}
% ============================================================================

\subsection{Activity 1: HITL Detection Hunt (15 minutes)}

\textbf{Setup:} Students work in pairs with their smartphones/laptops.

\textbf{Task:} Find and document 5 examples of systems asking for human input.

\paragraph{Examples to look for.}
\begin{itemize}
    \item ``Are you still there?'' prompts
    \item App permission requests
    \item ``Did you mean\ldots?'' suggestions
    \item ``Are you a robot?'' CAPTCHAs
    \item ``Is this your location?'' confirmations
    \item Two-factor authentication prompts
    \item Fraud alerts from banks
\end{itemize}

\paragraph{Analysis.} For each example, identify which of the five dimensions are visible:
\begin{description}
    \item[U] Can you tell the system is uncertain?
    \item[I] How does it ask?
    \item[T] When does it ask?
    \item[S] What are the stakes?
    \item[F] Does your response seem to affect future behavior?
\end{description}

\textbf{Debrief:} Which dimension is most often invisible to users?

\subsection{Activity 2: Framework Application Workshop (20 minutes)}

\textbf{Setup:} Small groups of 3--4 students.

\textbf{Task:} Apply the five-dimension framework to analyze one of these systems:
\begin{itemize}
    \item Google Photos face tagging
    \item Spotify music recommendations
    \item Amazon product recommendations
    \item Smart home device (Alexa, Google Home)
    \item Banking fraud detection
\end{itemize}

\textbf{Deliverable:} Completed framework analysis table + 1 recommended improvement.

\begin{table}[htbp]
\caption{Five-dimension framework analysis worksheet}
\label{tab:framework-worksheet}
\centering
\begin{tabular}{@{}p{3.5cm}p{2.5cm}p{3cm}p{3cm}@{}}
\toprule
\textbf{Dimension} & \textbf{Rating (1--5)} & \textbf{Evidence} & \textbf{Improvement} \\
\midrule
Uncertainty Detection & & & \\[8pt]
Intervention Design & & & \\[8pt]
Timing & & & \\[8pt]
Stakes Calibration & & & \\[8pt]
Feedback Integration & & & \\
\bottomrule
\end{tabular}
\end{table}

\subsection{Activity 3: Redesign Challenge (25 minutes)}

\textbf{Setup:} Individual or pairs.

\textbf{Task:} Redesign the Nest thermostat's Home/Away detection using the five-dimension framework.

\paragraph{Requirements.}
\begin{itemize}
    \item When should it ask vs.\ act autonomously?
    \item What should the question look like?
    \item How should stakes affect the threshold?
    \item How should it learn from user responses?
\end{itemize}

\textbf{Deliverable:} Interface mockup (sketch) + decision flowchart.

\subsection{Activity 4: Cost-Benefit Analysis (25 minutes)}

\textbf{Setup:} Individual work, then class discussion.

\textbf{Task:} Analyze the trade-offs for different intervention frequencies.

\paragraph{Scenario: Email spam filter.}
\begin{description}
    \item[Conservative] (asks about many emails): High accuracy, high annoyance
    \item[Aggressive] (rarely asks): Low accuracy, low annoyance
    \item[Balanced] (optimal threshold): Need to find sweet spot
\end{description}

\paragraph{Given.}
\begin{itemize}
    \item Cost of missed spam: \$2 (user time)
    \item Cost of blocked legitimate email: \$50 (business impact)
    \item Cost of asking user: \$0.10 (interruption)
    \item Spam detection accuracy: 95\% without human help, 99\% with help
\end{itemize}

\textbf{Exercise:} Calculate optimal intervention threshold.

\subsection{Activity 5: ``Try This'' Debrief (10 minutes)}

\textbf{Timing:} Beginning of second class session.

\textbf{Task:} Students share findings from the 24-hour HITL detection exercise assigned in Chapter~1.

\paragraph{Discussion Questions.}
\begin{itemize}
    \item How many HITL moments did you count?
    \item Which systems asked well? Which asked poorly?
    \item Did any systems that \emph{should} have asked, not ask?
    \item What patterns did you notice?
\end{itemize}

% ============================================================================
\section{Assessment Strategies}
% ============================================================================

\subsection{Formative Assessment (During Class)}

\subsubsection{Concept Checks}

\paragraph{Quick Polls.}
\begin{itemize}
    \item ``Which dimension of the framework does this represent?'' (show example)
    \item ``True or False: The best AI systems never need human help''
    \item ``Which failure mode is worse: asking too often, or never asking?''
\end{itemize}

\subsubsection{Exit Tickets}

\paragraph{3-2-1 Format.}
\begin{itemize}
    \item 3 examples of HITL systems you've used today
    \item 2 dimensions of the framework you can apply confidently
    \item 1 question you still have about the topic
\end{itemize}

\paragraph{Framework Application.}
``Apply the five-dimension framework to one system you used today'' (quick sketch).

\subsubsection{Think-Pair-Share}

\textbf{Prompt 1:} ``Explain to your partner why a voice assistant says `I don't understand' instead of just guessing what you meant.''

\textbf{Prompt 2:} ``Using the framework, explain why the Nest thermostat's Home/Away feature frustrates users.''

\subsection{Summative Assessment Options}

\subsubsection{Option 1: Framework Analysis Paper (750 words)}

\textbf{Prompt:} ``Choose a technology you use regularly that incorporates human-in-the-loop principles. Using the five-dimension framework, analyze how it handles uncertainty, evaluate its effectiveness, and propose one specific improvement.''

\paragraph{Rubric.}
\begin{itemize}
    \item \textbf{Framework Application (35\%)}: Correctly applies all five dimensions
    \item \textbf{Analysis Quality (25\%)}: Demonstrates understanding of uncertainty types and intervention strategies
    \item \textbf{Evidence (20\%)}: Uses specific examples and observations
    \item \textbf{Improvement Proposal (20\%)}: Proposes feasible, well-reasoned improvement
\end{itemize}

\subsubsection{Option 2: Design Challenge}

\textbf{Prompt:} ``Design a human-in-the-loop system for a specific application. Include uncertainty detection, intervention triggers, and user interface mockups.''

\paragraph{Deliverables.}
\begin{itemize}
    \item Five-dimension framework analysis
    \item System architecture diagram
    \item Interface mockups
    \item Decision flow chart
    \item 500-word design rationale
\end{itemize}

\paragraph{Assessment Focus.}
\begin{itemize}
    \item Understanding of HITL principles
    \item Practical application of five-dimension framework
    \item User-centered design thinking
    \item Technical feasibility
\end{itemize}

\subsubsection{Option 3: Case Study Presentation (10 minutes)}

\textbf{Assignment:} Research a real-world HITL system failure or success story (suggestions: Air Canada chatbot, GPS navigation failures, medical AI implementations).

\paragraph{Requirements.}
\begin{itemize}
    \item Technical explanation of what happened
    \item Analysis using the five-dimension framework
    \item Lessons learned for future systems
    \item Q\&A session with class
\end{itemize}

\subsubsection{Option 4: Technical Implementation (for advanced courses)}

\textbf{Task:} Implement uncertainty estimation for a simple classifier.

\paragraph{Components.}
\begin{itemize}
    \item Code implementing Monte Carlo Dropout or ensemble methods
    \item Analysis of uncertainty calibration
    \item Recommendation for human intervention thresholds
    \item Five-dimension framework analysis of your implementation
    \item Written reflection on challenges and insights
\end{itemize}

% ============================================================================
\section{Common Student Misconceptions \& How to Address Them}
% ============================================================================

\subsection{Misconception 1: ``Asking for help = system failure''}

\paragraph{Student thinking:} ``Good AI should never need human help.''

\paragraph{Correction strategy.}
\begin{itemize}
    \item Use medical analogy: ``Even experienced doctors get second opinions''
    \item Show Air Canada chatbot case: overconfidence caused legal liability
    \item Emphasize that recognition of limits is a form of intelligence
    \item Framework connection: ``Uncertainty Detection is a feature, not a bug''
\end{itemize}

\subsection{Misconception 2: ``All uncertainty is the same''}

\paragraph{Student thinking:} ``If the system isn't sure, just ask a human.''

\paragraph{Correction strategy.}
\begin{itemize}
    \item Clear examples of aleatoric vs.\ epistemic uncertainty
    \item Show how different uncertainty types require different responses
    \item Practice categorizing uncertainty sources using Chapter~1 examples
\end{itemize}

\subsection{Misconception 3: ``Humans are always better judges''}

\paragraph{Student thinking:} ``Human input always improves decisions.''

\paragraph{Correction strategy.}
\begin{itemize}
    \item Examples of human cognitive biases
    \item Discussion of expertise levels and context
    \item Alert fatigue research: too many questions = humans ignore them
    \item Framework connection: ``Stakes Calibration determines when human input is worth the cost''
\end{itemize}

\subsection{Misconception 4: ``This is just about AI systems''}

\paragraph{Student thinking:} ``HITL only applies to machine learning.''

\paragraph{Correction strategy.}
\begin{itemize}
    \item Broader examples: calculators asking ``are you sure?'', spell-check, etc.
    \item Historical perspective: Human-computer collaboration predates AI
    \item The Nest thermostat isn't ``AI'' but still exhibits HITL design choices
    \item Emphasize general principles of system design
\end{itemize}

\subsection{Misconception 5: ``The framework is just academic''}

\paragraph{Student thinking:} ``This is theoretical, not practical.''

\paragraph{Correction strategy.}
\begin{itemize}
    \item Show how the framework predicts which systems frustrate users (Nest)
    \item Use framework to diagnose real failures (Air Canada)
    \item Have students apply framework to systems they use daily
    \item Point to legal implications (companies liable for chatbot errors)
\end{itemize}

% ============================================================================
\section{Extension Activities \& Research Projects}
% ============================================================================

\subsection{Individual Research Topics}
\begin{enumerate}
    \item \textbf{Historical Analysis:} How did early computer systems handle uncertainty? Compare to modern approaches.
    \item \textbf{Cultural Differences:} How do uncertainty communication preferences vary across cultures?
    \item \textbf{Accessibility:} How should HITL systems be designed for users with disabilities?
    \item \textbf{Industry Deep Dive:} Analyze HITL implementation in a specific domain (healthcare, finance, transportation).
    \item \textbf{Legal Implications:} Research the Air Canada case and similar AI liability rulings.
\end{enumerate}

\subsection{Group Projects}
\begin{enumerate}
    \item \textbf{System Comparison:} Compare how different companies handle similar uncertainty scenarios (e.g., Google vs.\ Apple voice assistants) using the five-dimension framework.
    \item \textbf{User Study:} Design and conduct interviews about user experiences with HITL systems.
    \item \textbf{Design Challenge:} Propose HITL solutions for emerging technologies (VR/AR, IoT, etc.).
    \item \textbf{Failure Database:} Create a collection of HITL failures with framework analysis.
\end{enumerate}

\subsection{Advanced Technical Projects}
\begin{enumerate}
    \item \textbf{Implementation:} Build a working HITL system with uncertainty quantification.
    \item \textbf{Algorithm Development:} Improve uncertainty estimation for a specific type of model.
    \item \textbf{Evaluation Framework:} Develop new metrics for assessing HITL system effectiveness.
    \item \textbf{Intervention Optimization:} Implement and test different intervention timing strategies.
\end{enumerate}

% ============================================================================
\section{Technology Integration}
% ============================================================================

\subsection{Required Technology}
\begin{itemize}
    \item \textbf{Basic:} Computer/projector for demonstrations, internet access
    \item \textbf{Enhanced:} Student devices for hands-on activities
    \item \textbf{Advanced:} Access to programming environment for technical exercises
\end{itemize}

\subsection{Recommended Tools}

\paragraph{Demo Platforms.}
\begin{itemize}
    \item Teachable Machine (for quick ML demonstrations)
    \item Scratch for AI (visual programming examples)
    \item Online uncertainty visualization tools
\end{itemize}

\paragraph{Collaboration.}
\begin{itemize}
    \item Miro/Mural for system design activities and framework diagrams
    \item Padlet for collecting HITL examples
    \item Poll Everywhere for real-time class feedback
\end{itemize}

\subsection{Virtual/Hybrid Considerations}
\begin{itemize}
    \item \textbf{Screen sharing} for demonstrations works well
    \item \textbf{Breakout rooms} for small group activities
    \item \textbf{Shared documents} for collaborative note-taking
    \item \textbf{Recorded examples} can be paused and discussed
\end{itemize}

% ============================================================================
\section{Connections to Other Chapters}
% ============================================================================

\subsection{Backward Links}
\begin{itemize}
    \item \textbf{Prerequisites from Introduction:} Basic understanding of AI capabilities and limitations
    \item \textbf{Foundational concepts:} What is ``intelligence'' in artificial systems?
\end{itemize}

\subsection{Forward Links}
\begin{itemize}
    \item \textbf{Chapter 2:} How do we detect when systems are confused? (Deeper dive into Uncertainty Detection dimension)
    \item \textbf{Chapter 3:} How do machine learning systems actually learn from uncertainty?
    \item \textbf{Chapter 4:} Technical framework for HITL decision-making (formal treatment of the five dimensions)
    \item \textbf{Parts II--III:} Implementation and real-world examples build on these principles
\end{itemize}

\subsection{Cross-Curricular Connections}
\begin{itemize}
    \item \textbf{Psychology:} Human decision-making and cognitive biases (alert fatigue)
    \item \textbf{Philosophy:} Nature of intelligence and knowledge
    \item \textbf{Ethics:} Responsibility and transparency in AI systems (Air Canada case)
    \item \textbf{Business:} Cost-benefit analysis and user experience design
    \item \textbf{Statistics:} Probability, confidence intervals, uncertainty quantification
    \item \textbf{Law:} AI liability and corporate responsibility for chatbot errors
\end{itemize}

% ============================================================================
\section{Accessibility \& Inclusion Notes}
% ============================================================================

\subsection{Universal Design Principles}
\begin{itemize}
    \item \textbf{Multiple examples} to connect with diverse experiences (streaming, banking, smart home, medical)
    \item \textbf{Visual, auditory, and kinesthetic} learning opportunities
    \item \textbf{Scaffolded complexity} from concrete to abstract concepts
    \item \textbf{Cultural sensitivity} in examples and case studies
\end{itemize}

\subsection{Accommodation Strategies}
\begin{itemize}
    \item \textbf{Visual impairments:} Detailed verbal descriptions of screenshots and demos
    \item \textbf{Hearing impairments:} Visual demonstrations, written summaries
    \item \textbf{Learning differences:} Multiple explanation styles, extended time for activities
    \item \textbf{ESL considerations:} Clear technical vocabulary definitions, cultural context for examples
\end{itemize}

\subsection{Inclusive Examples}
\begin{itemize}
    \item \textbf{Global perspectives:} Air Canada case (Canada), examples from different countries
    \item \textbf{Diverse applications:} Healthcare, education, entertainment, accessibility tools
    \item \textbf{Varied expertise levels:} Examples that don't assume technical background
    \item \textbf{Multiple interaction styles:} Voice, touch, traditional interfaces
\end{itemize}

% ============================================================================
\section{Assessment Rubrics}
% ============================================================================

\subsection{Participation Rubric}

\begin{table}[htbp]
\caption{Participation rubric for discussion-based classes}
\label{tab:participation-rubric}
\centering
\begin{tabular}{@{}p{2.5cm}p{2.5cm}p{2.5cm}p{2.5cm}p{2.5cm}@{}}
\toprule
\textbf{Criterion} & \textbf{Excellent (4)} & \textbf{Good (3)} & \textbf{Satisfactory (2)} & \textbf{Needs Improvement (1)} \\
\midrule
Preparation & Clearly has done readings; references specific examples & Shows familiarity with material & Basic understanding evident & Little evidence of preparation \\
Contribution & Insightful comments that advance discussion & Relevant contributions, asks good questions & Participates when prompted & Minimal or off-topic contributions \\
Framework Use & Applies five-dimension framework accurately and creatively & Uses framework appropriately & Basic framework application & Struggles to apply framework \\
Analysis & Applies concepts creatively, sees connections & Uses chapter concepts appropriately & Basic application of ideas & Struggles to apply concepts \\
\bottomrule
\end{tabular}
\end{table}

\subsection{Framework Analysis Rubric}

\begin{table}[htbp]
\caption{Framework analysis rubric}
\label{tab:framework-rubric}
\centering
\begin{tabular}{@{}p{2.5cm}p{2.5cm}p{2.5cm}p{2.5cm}p{2.5cm}@{}}
\toprule
\textbf{Criterion} & \textbf{Excellent (4)} & \textbf{Good (3)} & \textbf{Satisfactory (2)} & \textbf{Needs Improvement (1)} \\
\midrule
Completeness & All five dimensions analyzed thoroughly & All five dimensions addressed & Most dimensions addressed & Missing multiple dimensions \\
Accuracy & Correct application of each dimension & Minor errors in application & Some misunderstandings & Fundamental misapplication \\
Evidence & Strong, specific evidence for each rating & Good evidence for most ratings & Some evidence provided & Assertions without evidence \\
Insight & Identifies non-obvious patterns and implications & Good analysis with some insight & Basic analysis & Superficial analysis \\
\bottomrule
\end{tabular}
\end{table}

\subsection{Technical Exercise Rubric}

\begin{table}[htbp]
\caption{Technical exercise rubric}
\label{tab:technical-rubric}
\centering
\begin{tabular}{@{}p{2.5cm}p{2.5cm}p{2.5cm}p{2.5cm}p{2.5cm}@{}}
\toprule
\textbf{Criterion} & \textbf{Excellent (4)} & \textbf{Good (3)} & \textbf{Satisfactory (2)} & \textbf{Needs Improvement (1)} \\
\midrule
Technical Accuracy & Code works correctly; proper implementation & Minor technical issues, mostly correct & Basic functionality, some errors & Significant technical problems \\
Conceptual Understanding & Clear grasp of uncertainty principles & Good understanding with minor gaps & Basic understanding & Fundamental misconceptions \\
Documentation & Clear explanations, well-commented code & Good documentation, mostly clear & Adequate documentation & Poor or missing documentation \\
Innovation & Creative solutions, goes beyond requirements & Some creative elements & Meets basic requirements & Below minimum requirements \\
\bottomrule
\end{tabular}
\end{table}

% ============================================================================
\section{Sample Lesson Plan}
% ============================================================================

\begin{table}[htbp]
\caption{Sample 50-minute lesson plan}
\label{tab:lesson-plan}
\centering
\begin{tabular}{@{}p{2.5cm}p{6cm}p{4cm}@{}}
\toprule
\textbf{Time} & \textbf{Activity} & \textbf{Materials} \\
\midrule
0:00--0:05 & Opening: Netflix demo + poll & Slide with Netflix screenshot \\
0:05--0:12 & Examples: Autocorrect, voice assistants, Nest & Slides with real examples \\
0:12--0:22 & Framework introduction + Netflix analysis & Framework diagram, whiteboard \\
0:22--0:37 & High-stakes case studies (GPS, Air Canada, Medical AI) & Case study slides \\
0:37--0:45 & Two failure modes discussion & Whiteboard diagram \\
0:45--0:50 & Wrap-up, assign ``Try This'' exercise & Chapter handout \\
\bottomrule
\end{tabular}
\end{table}

\bigskip
\begin{center}
\rule{0.5\textwidth}{0.4pt}
\end{center}
\bigskip

\noindent\textit{This teacher's manual provides everything needed to successfully teach Chapter~1 across different course levels and formats. The five-dimension framework provides a consistent analytical tool that students can apply throughout the course and beyond.}

\medskip
\noindent\textit{Next: Chapter~2 explores how systems detect uncertainty---diving deeper into the ``Uncertainty Detection'' dimension of the framework.}
